% =====================================================================
%  Banco complementar --- Indutores  (REVISADO)
%  Substitui a versao gerada por template (12 clones em N2, 12 em N3,
%  10 questoes conceituais indevidamente marcadas como N4).
%  Sem sobreposicao com o cap10, que ja traz: indutor ideal em CC (MC),
%  indutancia de solenoide N=500, rampa de corrente basica, paralelo
%  3H/9H, RL com L=1H e R=100 ohm, Leq serie/paralelo e a tabela
%  comparativa capacitor x indutor.
%  Distribuicao mantida: 6 N1 + 12 N2 + 12 N3 + 10 N4 = 40
% =====================================================================

\subsubsection*{Banco complementar --- Indutores}

\begin{atencao}[Organização do banco]
Três relações sustentam tudo: $v=L\dfrac{\mathrm{d}i}{\mathrm{d}t}$,
$W=\frac12LI^{2}$ e $\lambda=LI$. Delas seguem os dois fatos que governam o
comportamento transitório --- a corrente \textbf{não pode saltar} (isso exigiria
tensão infinita) e o indutor \textbf{se opõe à variação}, não à corrente em si.
O N3 trata de sobretensão, acoplamento, saturação e extração de parâmetros; o
N4 exige dedução da EDO, integração de energia e projeto com restrições
físicas.
\end{atencao}

\blocofixacao

\begin{questao}[1]{}
A tensão sobre um indutor vale $\SI{50}{\volt}$ quando a corrente varia à taxa
de $\SI{250}{\ampere\per\second}$. Determine a indutância.
\resposta{$L=\SI{0.2}{\henry}$}
\resolucao{Isolando $L$ na relação fundamental:
\[
v=L\frac{\mathrm{d}i}{\mathrm{d}t}
\;\Longrightarrow\;
L=\frac{v}{\mathrm{d}i/\mathrm{d}t}=\frac{50}{250}=\SI{0.2}{\henry}.
\]
O henry é, portanto, volt-segundo por ampère: $\SI{1}{\henry}$ é a indutância
que produz $\SI{1}{\volt}$ para uma variação de $\SI{1}{\ampere\per\second}$.}
\end{questao}

\begin{questao}[1]{}
Um indutor de $\SI{0.4}{\henry}$ conduz $\SI{5}{\ampere}$ constantes. Determine
a energia armazenada no campo magnético.
\resposta{$W=\SI{5}{\joule}$}
\resolucao{
\[
W=\frac12LI^{2}=\frac12(0{,}4)(5)^{2}=\SI{5}{\joule}.
\]
A dependência é \emph{quadrática} na corrente: dobrar $I$ quadruplica a
energia. Note que, embora a tensão sobre o indutor ideal seja nula em regime CC
(corrente constante), a energia armazenada não é --- ela ficou guardada no
campo durante a fase de crescimento da corrente.}
\end{questao}

\begin{questao}[1]{}
Sobre a corrente em um indutor ideal, é correto afirmar que ela:
\begin{alternativas}[2]
\item pode variar instantaneamente
\item não pode variar instantaneamente
\item é sempre constante
\item é sempre nula em corrente contínua
\end{alternativas}
\resposta{(b)}
\resolucao{Um salto de corrente significaria
$\mathrm{d}i/\mathrm{d}t\to\infty$ e, por $v=L\,\mathrm{d}i/\mathrm{d}t$, uma
tensão infinita --- fisicamente impossível. Logo a corrente no indutor é uma
função \textbf{contínua} do tempo: $i(0^{+})=i(0^{-})$.\\
Cuidado com (c) e (d): a corrente pode variar (e varia, nos transitórios), e em
CC ela é constante e em geral \emph{não nula}. O que é nula em CC é a
\emph{tensão}.}
\end{questao}

\begin{questao}[1]{}
Um indutor de $\SI{200}{\milli\henry}$ está em série com $\SI{50}{\ohm}$.
Determine a constante de tempo do circuito.
\resposta{$\tau=\SI{4}{\milli\second}$}
\resolucao{
\[
\tau=\frac{L}{R}=\frac{0{,}200}{50}=\SI{4e-3}{\second}=\SI{4}{\milli\second}.
\]
\textbf{Atenção à estrutura:} no circuito RL, $\tau=L/R$ --- resistência
\emph{maior} torna o transitório \emph{mais rápido}. No RC é o oposto
($\tau=RC$). Trocar as duas fórmulas é o erro mais comum entre os dois
elementos armazenadores.}
\end{questao}

\begin{questao}[1]{}
Determine o fluxo concatenado em um indutor de $\SI{0.25}{\henry}$ percorrido
por $\SI{8}{\ampere}$.
\resposta{$\lambda=\SI{2}{\weber}$}
\resolucao{
\[
\lambda=LI=0{,}25\cdot8=\SI{2}{\weber}.
\]
O fluxo concatenado $\lambda=N\Phi$ é o produto do fluxo por espira pelo número
de espiras. A indutância é justamente a constante de proporcionalidade entre
$\lambda$ e $I$ --- é essa a definição primária de $L$, da qual
$v=L\,\mathrm{d}i/\mathrm{d}t$ decorre pela lei de Faraday.}
\end{questao}

\begin{questao}[1]{}
Aplicam-se $\SI{30}{\volt}$ a um indutor de $\SI{0.15}{\henry}$ inicialmente
sem corrente. Determine a taxa inicial de variação da corrente.
\resposta{$\mathrm{d}i/\mathrm{d}t=\SI{200}{\ampere\per\second}$}
\resolucao{
\[
\frac{\mathrm{d}i}{\mathrm{d}t}=\frac{v}{L}=\frac{30}{0{,}15}
=\SI{200}{\ampere\per\second}.
\]
No instante inicial a corrente é nula, mas sua \emph{taxa} é máxima. Se a
tensão se mantivesse constante e não houvesse resistência, a corrente cresceria
indefinidamente em rampa --- é a resistência do circuito que estabelece o
patamar final.}
\end{questao}

\blocoaplicacao

\begin{questao}[2]{}
Um indutor de $\SI{0.5}{\henry}$ conduz $\SI{3}{\ampere}$.
\begin{itensq}
\item Determine a energia armazenada.
\item Determine a tensão necessária para levar a corrente a zero
      \emph{linearmente} em $\SI{50}{\milli\second}$.
\end{itensq}
\resposta{(a) $\SI{2.25}{\joule}$; (b) $\SI{30}{\volt}$}
\resolucao{(a) $W=\frac12(0{,}5)(9)=\SI{2.25}{\joule}$.\\
(b) A taxa é $\mathrm{d}i/\mathrm{d}t=(0-3)/0{,}050
=\SI{-60}{\ampere\per\second}$, e
\[
|v|=L\left|\frac{\mathrm{d}i}{\mathrm{d}t}\right|=0{,}5\cdot60=\SI{30}{\volt}.
\]
O sinal negativo indica \textbf{inversão de polaridade}: o indutor passa a
funcionar como fonte, devolvendo ao circuito os $\SI{2.25}{\joule}$
armazenados. É esse comportamento que produz as sobretensões de abertura.}
\end{questao}

\begin{questao}[2]{}
A corrente em um indutor de $\SI{100}{\milli\henry}$ cai linearmente de
$\SI{4}{\ampere}$ a zero em $\SI{20}{\milli\second}$. Determine a tensão
induzida e sua polaridade.
\resposta{$\SI{20}{\volt}$, com polaridade \emph{invertida} em relação à fase
de crescimento}
\resolucao{
\[
|v|=L\left|\frac{\Delta i}{\Delta t}\right|
=0{,}100\cdot\frac{4}{0{,}020}=\SI{20}{\volt}.
\]
Durante o \emph{crescimento} da corrente, o indutor absorve energia e sua
tensão se opõe ao aumento. Durante o \emph{decrescimento}, ele devolve energia
e a tensão inverte, tentando sustentar a corrente.\\
Em ambos os casos vale a lei de Lenz: o indutor se opõe à \emph{variação}, não
ao sentido da corrente.}
\end{questao}

\begin{questao}[2]{}
Três indutores não acoplados de $\SI{6}{\milli\henry}$,
$\SI{12}{\milli\henry}$ e $\SI{12}{\milli\henry}$ são associados. Determine
$L_{eq}$ em série e em paralelo.
\resposta{Série: $\SI{30}{\milli\henry}$; paralelo: $\SI{3}{\milli\henry}$}
\resolucao{\textbf{Série:} $L_{eq}=6+12+12=\SI{30}{\milli\henry}$.\\
\textbf{Paralelo:}
\[
\frac{1}{L_{eq}}=\frac16+\frac1{12}+\frac1{12}=\frac{2+1+1}{12}=\frac13
\;\Longrightarrow\;
L_{eq}=\SI{3}{\milli\henry}.
\]
As regras são \emph{idênticas} às dos resistores --- e opostas às dos
capacitores. A ressalva é a hipótese de \textbf{não acoplamento}: se houver
fluxo mútuo, as fórmulas mudam (ver bloco de aprofundamento).}
\end{questao}

\begin{questao}[2]{}
Um circuito RL série tem $L=\SI{0.6}{\henry}$, $R=\SI{30}{\ohm}$ e fonte de
$\SI{24}{\volt}$, com a chave fechada em $t=0$.
\begin{itensq}
\item Determine $\tau$ e a corrente final.
\item Determine a tensão sobre o indutor em $t=0^{+}$.
\end{itensq}
\resposta{(a) $\tau=\SI{20}{\milli\second}$, $I_f=\SI{0.8}{\ampere}$;
(b) $\SI{24}{\volt}$}
\resolucao{(a) $\tau=L/R=0{,}6/30=\SI{20}{\milli\second}$;
$I_f=V/R=24/30=\SI{0.8}{\ampere}$.\\
(b) Em $t=0^{+}$ a corrente ainda é nula (continuidade), logo o resistor não
tem queda e \textbf{toda} a tensão da fonte aparece sobre o indutor:
$v_L(0^{+})=\SI{24}{\volt}$.\\
\textbf{Os dois extremos:} em $t=0^{+}$ o indutor se comporta como circuito
\emph{aberto} (corrente nula); em $t\to\infty$, como \emph{curto}
(tensão nula). Reconhecer esses dois limites resolve a maioria dos problemas
sem integrar nada.}
\end{questao}

\begin{questao}[2]{}
No circuito anterior, determine a corrente em $t=\tau$, $t=2\tau$ e
$t=3\tau$.
\resposta{$\SI{0.506}{\ampere}$, $\SI{0.692}{\ampere}$ e
$\SI{0.760}{\ampere}$}
\resolucao{Com $i(t)=I_f\left(1-e^{-t/\tau}\right)$ e
$I_f=\SI{0.8}{\ampere}$:
\[
i(\tau)=0{,}8(1-e^{-1})=0{,}8(0{,}632)=\SI{0.506}{\ampere},
\]
\[
i(2\tau)=0{,}8(0{,}865)=\SI{0.692}{\ampere},
\qquad
i(3\tau)=0{,}8(0{,}950)=\SI{0.760}{\ampere}.
\]
As frações $63{,}2\%$, $86{,}5\%$ e $95{,}0\%$ são universais --- valem para
qualquer RL ou RC. Após $5\tau$ atinge-se $99{,}3\%$, e por convenção o
transitório é dado por encerrado.}
\end{questao}

\begin{questao}[2]{}
Aplica-se a um indutor de $\SI{5}{\milli\henry}$ uma onda quadrada de tensão
que alterna entre $+\SI{10}{\volt}$ e $-\SI{10}{\volt}$, com meio período de
$\SI{1}{\milli\second}$. Determine a variação de corrente pico a pico em
regime.
\resposta{$\Delta i=\SI{2}{\ampere}$}
\resolucao{Durante o semiciclo positivo, a tensão constante produz rampa de
corrente:
\[
\Delta i=\frac{v\,\Delta t}{L}=\frac{10\cdot0{,}001}{0{,}005}=\SI{2}{\ampere}.
\]
No semiciclo negativo a corrente desce pela mesma quantidade, e o regime é
simétrico.\\
\textbf{Resultado geral:} tensão quadrada no indutor produz corrente
\textbf{triangular} --- o indutor \emph{integra} a tensão. O capacitor faz o
oposto: corrente quadrada nele produz tensão triangular.}
\end{questao}

\begin{questao}[2]{}
Um indutor \emph{real} de $\SI{0.5}{\henry}$ possui resistência de enrolamento
$r=\SI{8}{\ohm}$ e é ligado a uma fonte CC de $\SI{24}{\volt}$.
\begin{itensq}
\item Determine a corrente em regime permanente.
\item Determine a energia armazenada e a potência dissipada.
\end{itensq}
\resposta{(a) $\SI{3}{\ampere}$; (b) $W=\SI{2.25}{\joule}$ e
$P=\SI{72}{\watt}$}
\resolucao{(a) Em regime, $\mathrm{d}i/\mathrm{d}t=0$ e o indutor ideal não cai
tensão: só $r$ limita a corrente.
\[
I=\frac{24}{8}=\SI{3}{\ampere}.
\]
(b) $W=\frac12(0{,}5)(9)=\SI{2.25}{\joule}$;
$P=rI^{2}=8(9)=\SI{72}{\watt}$.\\
\textbf{Contraste essencial:} a energia é \emph{armazenada} e recuperável; a
potência é \emph{dissipada} continuamente e perdida. Manter esses
$\SI{2.25}{\joule}$ no campo custa $\SI{72}{\watt}$ permanentes --- razão pela
qual eletroímãs de grande porte usam supercondutores, em que $r=0$ e a
manutenção do campo é gratuita.}
\end{questao}

\begin{questao}[2]{}
Dois indutores armazenam a mesma energia: o primeiro tem
$L_1=\SI{2}{\henry}$ com $I_1=\SI{1}{\ampere}$. Se
$L_2=\SI{0.5}{\henry}$, determine $I_2$.
\resposta{$I_2=\SI{2}{\ampere}$}
\resolucao{
\[
\frac12L_1I_1^{2}=\frac12L_2I_2^{2}
\;\Longrightarrow\;
I_2=I_1\sqrt{\frac{L_1}{L_2}}=1\cdot\sqrt{4}=\SI{2}{\ampere}.
\]
Ambos armazenam $\SI{1}{\joule}$. Reduzir $L$ por um fator $4$ exige
\emph{dobrar} a corrente --- e, como as perdas vão com $I^{2}$, o indutor menor
custa quatro vezes mais em dissipação para armazenar o mesmo tanto. É o
compromisso central no projeto de indutores de potência.}
\end{questao}

\begin{questao}[2]{}
Um solenoide de $N=200$ espiras tem indutância $L$. Determine a nova
indutância se o número de espiras for dobrado, mantidos o comprimento e a
seção.
\resposta{$4L$}
\resolucao{Para um solenoide longo,
\[
L=\frac{\mu N^{2}A}{\ell},
\]
de modo que $L\propto N^{2}$. Dobrar $N$ quadruplica $L$.\\
\textbf{Por que o quadrado:} dobrar as espiras dobra o campo produzido por
uma dada corrente \emph{e} dobra o número de espiras que concatenam esse campo.
Os dois efeitos se multiplicam.\\
\textbf{Ressalva prática:} com o mesmo espaço de bobina, dobrar $N$ exige fio
de seção menor, o que aumenta $r$ e as perdas. Na prática, o ganho de $L$ vem
acompanhado de penalidade em corrente admissível.}
\end{questao}

\begin{questao}[2]{}
Um indutor com núcleo de ar tem $L_0=\SI{2}{\milli\henry}$. Introduz-se um
núcleo ferromagnético de permeabilidade relativa $\mu_r=500$. Determine a nova
indutância e comente as limitações.
\resposta{$L=\SI{1}{\henry}$}
\resolucao{Como $L\propto\mu=\mu_r\mu_0$:
\[
L=\mu_rL_0=500\cdot2=\SI{1000}{\milli\henry}=\SI{1}{\henry}.
\]
\textbf{Três limitações do núcleo.} \emph{(i)} \textbf{Saturação:} acima de
certa corrente, $\mu_r$ despenca e $L$ colapsa. \emph{(ii)} \textbf{Perdas:}
histerese e correntes de Foucault dissipam energia proporcionalmente à
frequência. \emph{(iii)} \textbf{Não linearidade:} $L$ deixa de ser constante,
e as fórmulas lineares passam a ser aproximações válidas apenas em torno de um
ponto de operação.\\
Núcleos de ar têm $L$ minúsculo, mas são perfeitamente lineares e sem perdas
magnéticas --- por isso dominam aplicações de radiofrequência.}
\end{questao}

\begin{questao}[2]{}
No circuito RL com $\tau=\SI{20}{\milli\second}$ e $I_f=\SI{0.8}{\ampere}$,
determine o instante em que a corrente atinge $\SI{0.5}{\ampere}$.
\resposta{$t\approx\SI{19.6}{\milli\second}$}
\resolucao{De $i=I_f(1-e^{-t/\tau})$, isole a exponencial:
\[
e^{-t/\tau}=1-\frac{0{,}5}{0{,}8}=0{,}375
\;\Longrightarrow\;
t=\tau\ln\!\left(\frac{1}{0{,}375}\right)=0{,}020\cdot0{,}981
=\SI{19.6}{\milli\second}.
\]
Ou seja, praticamente $1\tau$ --- coerente com o fato de que
$\SI{0.5}{\ampere}$ corresponde a $62{,}5\%$ de $I_f$, muito próximo dos
$63{,}2\%$ atingidos em exatamente $\tau$.}
\end{questao}

\begin{questao}[2]{}
Um indutor de $\SI{0.4}{\henry}$ conduz $\SI{2}{\ampere}$ quando a fonte é
substituída por um curto através de $R=\SI{20}{\ohm}$. Determine $\tau$ e a
corrente após $\SI{40}{\milli\second}$.
\resposta{$\tau=\SI{20}{\milli\second}$; $i=\SI{0.27}{\ampere}$}
\resolucao{$\tau=L/R=0{,}4/20=\SI{20}{\milli\second}$.\\
Na descarga, a corrente decai exponencialmente a partir do valor inicial:
\[
i(t)=I_0e^{-t/\tau}=2\,e^{-40/20}=2e^{-2}=2(0{,}135)=\SI{0.27}{\ampere}.
\]
\textbf{Continuidade:} no instante da comutação a corrente permanece
$\SI{2}{\ampere}$ --- ela não pode saltar. O que muda instantaneamente é a
\emph{tensão}, que assume $-RI_0=\SI{-40}{\volt}$ para manter essa corrente
circulando pelo resistor.}
\end{questao}

\blocoaprofundamento

\begin{questao}[3]{}
Um indutor de $\SI{0.2}{\henry}$ conduz $\SI{5}{\ampere}$ e é descarregado
através de um resistor.
\begin{itensq}
\item Determine a energia total dissipada no resistor.
\item Explique por que ela não depende do valor de $R$.
\item Determine como $R$ afeta o processo.
\end{itensq}
\resposta{(a) $\SI{2.5}{\joule}$; (b) conservação de energia; (c) $R$ altera a
\emph{duração} e a \emph{potência}, não a energia total}
\resolucao{(a) Toda a energia armazenada acaba no resistor:
\[
W=\frac12LI_0^{2}=\frac12(0{,}2)(25)=\SI{2.5}{\joule}.
\]
(b) O indutor ideal não dissipa e o circuito é fechado: a energia magnética não
tem outro destino. A independência em relação a $R$ é consequência direta da
conservação --- e é confirmada por integração no bloco de desafio.\\
(c) $R$ governa $\tau=L/R$ e, portanto, o \emph{ritmo}:
\begin{itemize}
\item $R$ pequeno $\Rightarrow$ $\tau$ grande: descarga lenta, potência baixa e
prolongada.
\item $R$ grande $\Rightarrow$ $\tau$ pequeno: descarga rápida, com pico de
potência $P_0=RI_0^{2}$ muito elevado.
\end{itemize}
\textbf{Exemplo:} com $R=\SI{10}{\ohm}$, $P_0=\SI{250}{\watt}$ por
$\SI{20}{\milli\second}$; com $R=\SI{1000}{\ohm}$, $P_0=\SI{25}{\kilo\watt}$
por $\SI{0.2}{\milli\second}$. Mesma energia, potências que diferem em cem
vezes.}
\end{questao}

\begin{questao}[3]{}
Um indutor de $\SI{0.1}{\henry}$ conduz $\SI{2}{\ampere}$ quando uma chave
mecânica é aberta, interrompendo a corrente em cerca de
$\SI{1}{\micro\second}$.
\begin{itensq}
\item Estime a tensão induzida.
\item Explique por que esse valor não se observa na prática.
\item Indique a consequência real.
\end{itensq}
\resposta{(a) $\SI{200}{\kilo\volt}$ (estimativa idealizada);
(b) o ar rompe muito antes; (c) formação de arco elétrico nos contatos}
\resolucao{(a)
\[
|v|=L\frac{\Delta i}{\Delta t}=0{,}1\cdot\frac{2}{10^{-6}}
=\pot{2}{5}\,\si{\volt}=\SI{200}{\kilo\volt}.
\]
(b) A rigidez dielétrica do ar é da ordem de
$\SI{3}{\kilo\volt\per\milli\meter}$. Com contatos separados por décimos de
milímetro, o ar ioniza a alguns \emph{quilovolts} --- muito antes dos
$\SI{200}{\kilo\volt}$. O cálculo é uma \textbf{estimativa por absurdo}: ele
mostra que a hipótese "a corrente cessa em $\SI{1}{\micro\second}$" é
insustentável.\\
(c) Forma-se um \textbf{arco elétrico}, que mantém a corrente circulando
enquanto o campo se esvazia. O arco limita a tensão a algumas centenas de
volts, mas corrói os contatos, gera interferência eletromagnética e reduz
drasticamente a vida da chave.\\
\textbf{Leitura correta:} não é o indutor que "gera" $\SI{200}{\kilo\volt}$ ---
é a impossibilidade de interromper a corrente que \emph{força} o circuito a
encontrar um caminho, e ele o encontra no ar.}
\end{questao}

\begin{questao}[3]{}
Um diodo de roda livre é colocado em antiparalelo com uma bobina de
$\SI{0.1}{\henry}$ e $r=\SI{5}{\ohm}$, que conduz $\SI{2}{\ampere}$.
\begin{itensq}
\item Explique o funcionamento no instante da abertura.
\item Determine a tensão máxima sobre a chave.
\item Determine a constante de tempo de extinção.
\end{itensq}
\resposta{(b) $\approx\SI{0.7}{\volt}$ acima da alimentação;
(c) $\tau=\SI{20}{\milli\second}$}
\resolucao{(a) Com a chave aberta, a corrente da bobina precisa de um caminho.
O diodo, que estava reversamente polarizado, entra em condução e fecha uma
malha bobina--diodo. A corrente continua circulando e decai exponencialmente,
dissipando-se em $r$ e no diodo --- sem arco.\\
(b) A tensão sobre a bobina fica grampeada em $-V_D\approx\SI{-0.7}{\volt}$, de
modo que a chave suporta apenas a tensão de alimentação acrescida de
$\SI{0.7}{\volt}$. Compare com os quilovolts do caso anterior.\\
(c) $\tau=L/r=0{,}1/5=\SI{20}{\milli\second}$.\\
\textbf{Custo do método:} a extinção fica \emph{lenta}. Em um relé, isso
prolonga o tempo de desarme; em acionamentos rápidos, prefere-se um diodo em
série com um resistor (ou um Zener), que eleva a tensão de grampeamento e
acelera a queda --- trocando proteção máxima por velocidade.}
\end{questao}

\begin{questao}[3]{}
Dois indutores acoplados, $L_1=\SI{4}{\milli\henry}$ e
$L_2=\SI{9}{\milli\henry}$, têm indutância mútua $M=\SI{3}{\milli\henry}$.
\begin{itensq}
\item Determine $L_{eq}$ em série nas duas orientações possíveis.
\item Determine o coeficiente de acoplamento $k$.
\item Explique como identificar experimentalmente a orientação.
\end{itensq}
\resposta{(a) $\SI{19}{\milli\henry}$ e $\SI{7}{\milli\henry}$; (b) $k=0{,}5$}
\resolucao{(a) Em série,
\[
L_{eq}=L_1+L_2\pm2M
\;\Longrightarrow\;
4+9+6=\SI{19}{\milli\henry}
\ \text{ou}\ 4+9-6=\SI{7}{\milli\henry}.
\]
O sinal $+$ vale para \textbf{acoplamento aditivo} (os fluxos se reforçam; a
corrente entra pelos dois pontos marcados); o sinal $-$, para
\textbf{subtrativo}.\\
(b)
\[
k=\frac{M}{\sqrt{L_1L_2}}=\frac{3}{\sqrt{36}}=0{,}5.
\]
Como $0\le k\le1$, este é um acoplamento médio --- típico de bobinas
próximas com núcleo de ar. Transformadores de potência atingem $k>0{,}99$.\\
(c) \textbf{Método experimental:} meça $L_{eq}$ nas duas ligações possíveis. A
\emph{maior} corresponde ao acoplamento aditivo e identifica os pontos.
Além disso, $M$ sai de graça:
\[
M=\frac{L_{aditivo}-L_{subtrativo}}{4}=\frac{19-7}{4}=\SI{3}{\milli\henry}.
\]
É assim que se determina $M$ sem acesso à geometria interna das bobinas.}
\end{questao}

\begin{questao}[3]{}
Uma fonte de $\SI{12}{\volt}$ alimenta, através de $R_1=\SI{20}{\ohm}$, um nó
do qual saem $R_3=\SI{60}{\ohm}$ ao terra e um ramo com
$R_2=\SI{30}{\ohm}$ em série com $L=\SI{0.9}{\henry}$.
\begin{itensq}
\item Determine a resistência de Thévenin vista pelo indutor.
\item Determine $\tau$ e a corrente final no indutor.
\end{itensq}
\resposta{(a) $R_{Th}=\SI{45}{\ohm}$; (b) $\tau=\SI{20}{\milli\second}$ e
$I_f=\SI{0.2}{\ampere}$}
\resolucao{(a) Zerando a fonte (curto), o indutor vê $R_2$ em série com
$R_1\parallel R_3$:
\[
R_{Th}=30+\frac{20\cdot60}{80}=30+15=\SI{45}{\ohm}.
\]
(b) $\tau=L/R_{Th}=0{,}9/45=\SI{20}{\milli\second}$.\\
A tensão de Thévenin sai com o indutor removido (sem corrente em $R_2$, logo
sem queda nele):
\[
V_{Th}=12\cdot\frac{60}{20+60}=\SI{9}{\volt}
\;\Longrightarrow\;
I_f=\frac{9}{45}=\SI{0.2}{\ampere}.
\]
\textbf{Método geral:} em qualquer circuito com \emph{um} elemento armazenador,
substitua todo o restante pelo equivalente de Thévenin visto por ele. O
problema se reduz a um RL série, com $\tau=L/R_{Th}$ e
$I_f=V_{Th}/R_{Th}$ --- e não é preciso resolver a rede inteira em regime
transitório.}
\end{questao}

\begin{questao}[3]{}
Um indutor de $\SI{0.5}{\henry}$ já conduz $\SI{1}{\ampere}$ quando é ligado a
uma fonte que estabeleceria $\SI{3}{\ampere}$ em regime, com
$R=\SI{25}{\ohm}$.
\begin{itensq}
\item Escreva $i(t)$.
\item Determine a corrente após $\SI{20}{\milli\second}$.
\item Determine a tensão inicial sobre o indutor.
\end{itensq}
\resposta{(a) $i(t)=3-2e^{-t/\tau}$ com $\tau=\SI{20}{\milli\second}$;
(b) $\SI{2.26}{\ampere}$; (c) $\SI{50}{\volt}$}
\resolucao{(a) A forma geral do transitório de primeira ordem é
\[
i(t)=I_f+\left(I_0-I_f\right)e^{-t/\tau},
\]
com $I_0=\SI{1}{\ampere}$, $I_f=\SI{3}{\ampere}$ e
$\tau=0{,}5/25=\SI{20}{\milli\second}$:
\[
i(t)=3-2e^{-t/0{,}020}.
\]
(b) $i(0{,}020)=3-2e^{-1}=3-0{,}736=\SI{2.26}{\ampere}$.\\
(c) Em $t=0^{+}$, $i=\SI{1}{\ampere}$ e o resistor cai
$25(1)=\SI{25}{\volt}$; a fonte vale $RI_f=25(3)=\SI{75}{\volt}$. Logo
\[
v_L(0^{+})=75-25=\SI{50}{\volt}.
\]
\textbf{Generalidade:} a fórmula de (a) cobre carga, descarga e condições
iniciais quaisquer --- basta identificar $I_0$, $I_f$ e $\tau$. Não é preciso
memorizar três expressões diferentes.}
\end{questao}

\begin{questao}[3]{}
Em um circuito RL, mediu-se corrente final de $\SI{2}{\ampere}$ com fonte de
$\SI{50}{\volt}$, e verificou-se que a corrente atingiu $\SI{1.264}{\ampere}$
em $\SI{8}{\milli\second}$. Determine $R$ e $L$.
\resposta{$R=\SI{25}{\ohm}$; $L=\SI{0.2}{\henry}$}
\resolucao{\textbf{Resistência:} do regime permanente,
$R=V/I_f=50/2=\SI{25}{\ohm}$.\\
\textbf{Constante de tempo:} note que
$1{,}264/2=0{,}632$ --- exatamente a fração característica de $1\tau$. Logo
$\tau=\SI{8}{\milli\second}$, e
\[
L=R\tau=25(0{,}008)=\SI{0.2}{\henry}.
\]
\textbf{Se a fração não fosse reconhecível}, o caminho geral seria
$\tau=-t/\ln(1-i/I_f)$.\\
\textbf{Método:} o regime permanente entrega $R$; o transitório entrega $\tau$
e, por consequência, $L$. Duas medidas de naturezas diferentes são necessárias,
porque nenhuma delas isola $L$ sozinha.}
\end{questao}

\begin{questao}[3]{}
Analise o sinal da potência instantânea $p=vi$ em um indutor durante um ciclo
completo de carga e descarga, e relacione-o com a energia armazenada.
\resposta{$p>0$ na carga (absorve), $p<0$ na descarga (devolve); a integral em
um ciclo completo é nula}
\resolucao{\textbf{Carga:} $i>0$ e crescente $\Rightarrow$
$v=L\,\mathrm{d}i/\mathrm{d}t>0$ $\Rightarrow$ $p>0$: o indutor \emph{absorve}
energia da fonte e a armazena no campo.\\
\textbf{Descarga:} $i>0$ mas decrescente $\Rightarrow$ $v<0$ $\Rightarrow$
$p<0$: o indutor \emph{fornece} energia ao circuito.\\
\textbf{Integral:} como
\[
\int p\,\mathrm{d}t=\int Li\,\frac{\mathrm{d}i}{\mathrm{d}t}\,\mathrm{d}t
=\int Li\,\mathrm{d}i=\frac12Li^{2},
\]
a energia depende \emph{apenas} da corrente instantânea, não do caminho
percorrido. Se o ciclo retorna à mesma corrente, a integral é nula: o indutor
ideal não consome energia líquida.\\
\textbf{Consequência:} indutor e capacitor são elementos \textbf{reativos} ---
trocam energia com o circuito sem consumi-la. Só a resistência dissipa. Essa
distinção é a base de todo o tratamento de potência ativa e reativa em corrente
alternada.}
\end{questao}

\begin{questao}[3]{}
Determine a densidade de energia magnética em um núcleo onde
$B=\SI{1.2}{\tesla}$, supondo $\mu\approx\mu_0$ no entreferro.
\begin{itensq}
\item Calcule $u$.
\item Compare com a densidade de energia elétrica em um campo de
      $\SI{3}{\mega\volt\per\meter}$.
\end{itensq}
\dados{$\mu_0=4\pi\times10^{-7}\,\si{\henry\per\meter}$;
$\varepsilon_0=\pot{8.85}{-12}\,\si{\farad\per\meter}$}
\resposta{(a) $u\approx\SI{573}{\kilo\joule\per\meter\cubed}$;
(b) $\approx\SI{40}{\joule\per\meter\cubed}$}
\resolucao{(a)
\[
u_B=\frac{B^{2}}{2\mu_0}=\frac{(1{,}2)^{2}}{2(4\pi\times10^{-7})}
=\frac{1{,}44}{\pot{2.513}{-6}}
\approx\pot{5.73}{5}\,\si{\joule\per\meter\cubed}.
\]
(b)
\[
u_E=\frac12\varepsilon_0E^{2}
=\frac12(\pot{8.85}{-12})(3\times10^{6})^{2}
\approx\SI{40}{\joule\per\meter\cubed}.
\]
\textbf{Razão: cerca de $14\,000$ vezes.} O campo elétrico usado é o limite de
ruptura do ar, e o magnético é típico de aço-silício --- ou seja, comparam-se
dois valores \emph{praticáveis}.\\
\textbf{Consequência tecnológica:} armazenar energia magneticamente é muito
mais denso do que eletrostaticamente. É por isso que transformadores e motores
operam com campos magnéticos, e por que capacitores só competem em energia
quando se usa dielétrico de altíssima constante e espessura nanométrica ---
como nos supercapacitores.}
\end{questao}

\begin{questao}[3]{}
Duas bobinas têm $L_1=\SI{50}{\milli\henry}$ e $L_2=\SI{200}{\milli\henry}$.
\begin{itensq}
\item Determine o valor máximo possível de $M$.
\item Se $M=\SI{80}{\milli\henry}$, determine $k$.
\item Interprete fisicamente $k=1$ e $k=0$.
\end{itensq}
\resposta{(a) $M_{\max}=\SI{100}{\milli\henry}$; (b) $k=0{,}8$}
\resolucao{(a) Como $k\le1$,
\[
M_{\max}=\sqrt{L_1L_2}=\sqrt{50\cdot200}=\sqrt{10\,000}
=\SI{100}{\milli\henry}.
\]
(b) $k=80/100=0{,}8$ --- acoplamento forte, típico de bobinas sobre um mesmo
núcleo ferromagnético com pequeno entreferro.\\
(c) $k=1$ significa \textbf{acoplamento perfeito}: \emph{todo} o fluxo produzido
por uma bobina concatena a outra, sem dispersão. É o transformador ideal ---
inatingível, mas aproximado em núcleos toroidais fechados de alta
permeabilidade.\\
$k=0$ significa \textbf{desacoplamento total}: nenhuma influência mútua. Obtém-se
afastando as bobinas ou dispondo seus eixos perpendicularmente, de modo que o
fluxo de uma atravesse a outra tangencialmente. É a razão de indutores vizinhos
em placas serem montados em orientações ortogonais.}
\end{questao}

\begin{questao}[3]{}
Um indutor com núcleo apresenta $L=\SI{100}{\micro\henry}$ até
$\SI{3}{\ampere}$; acima disso, a saturação reduz $L$ a
$\SI{20}{\micro\henry}$. Ele opera com $\SI{12}{\volt}$ aplicados durante
$\SI{10}{\micro\second}$ a partir de $\SI{2}{\ampere}$.
\begin{itensq}
\item Determine a corrente ao final do pulso, ignorando a saturação.
\item Refaça considerando a saturação.
\item Comente a consequência prática.
\end{itensq}
\resposta{(a) $\SI{3.2}{\ampere}$; (b) $\SI{5}{\ampere}$; (c) risco de
destruição do interruptor}
\resolucao{(a) Com $L$ constante,
\[
\Delta i=\frac{v\,\Delta t}{L}=\frac{12\cdot10^{-5}}{10^{-4}}
=\SI{1.2}{\ampere}
\;\Longrightarrow\;
i_{final}=2+1{,}2=\SI{3.2}{\ampere}.
\]
(b) A corrente sobe de $2$ a $\SI{3}{\ampere}$ com
$L=\SI{100}{\micro\henry}$, o que consome
\[
\Delta t_1=\frac{L\,\Delta i}{v}=\frac{10^{-4}\cdot1}{12}
=\SI{8.33}{\micro\second}.
\]
Restam $\SI{1.67}{\micro\second}$, agora com $L=\SI{20}{\micro\henry}$:
\[
\Delta i_2=\frac{12\cdot\pot{1.67}{-6}}{\pot{2}{-5}}=\SI{1}{\ampere}
\;\Longrightarrow\;
i_{final}=\SI{5}{\ampere}.
\]
(c) A corrente real é $56\%$ maior que a prevista pelo modelo linear, e a
subida final é \emph{cinco vezes} mais íngreme. Em um conversor chaveado, isso
significa que o transistor pode ver o dobro da corrente esperada em uma fração
do tempo --- destruição por sobrecorrente antes que qualquer proteção lenta
atue.\\
\textbf{Regra de projeto:} dimensione o indutor para que a corrente de pico
fique confortavelmente abaixo do joelho de saturação, e verifique
$L$ \emph{na corrente de operação}, nunca no valor de catálogo medido a
corrente nula.}
\end{questao}

\begin{questao}[3]{}
Uma fonte de \textbf{corrente} de $\SI{2}{\ampere}$ alimenta um indutor de
$\SI{0.5}{\henry}$ em paralelo com $R=\SI{10}{\ohm}$, em regime permanente.
\begin{itensq}
\item Determine a corrente no indutor e no resistor.
\item A fonte é subitamente desligada (aberta). Determine a tensão inicial
      sobre o resistor.
\item Determine o tempo de extinção.
\end{itensq}
\resposta{(a) $\SI{2}{\ampere}$ no indutor, $\SI{0}{\ampere}$ no resistor;
(b) $\SI{-20}{\volt}$; (c) $\tau=\SI{50}{\milli\second}$}
\resolucao{(a) Em regime CC o indutor ideal é um \textbf{curto}: ele desvia
toda a corrente, e o resistor em paralelo fica com tensão nula --- logo
corrente nula. Todos os $\SI{2}{\ampere}$ passam por $L$.\\
(b) Ao abrir a fonte, a corrente do indutor \emph{não pode saltar}: continua
$\SI{2}{\ampere}$, e o único caminho disponível é o resistor. A corrente agora
o percorre em sentido oposto ao anterior:
\[
v_R=-RI_0=-10(2)=\SI{-20}{\volt}.
\]
(c) $\tau=L/R=0{,}5/10=\SI{50}{\milli\second}$; a extinção é praticamente
completa em $5\tau=\SI{250}{\milli\second}$.\\
\textbf{Observação:} o resistor, que parecia inútil em regime (corrente nula),
é justamente o que evita a sobretensão destrutiva na abertura. Ele é um
\emph{caminho de escape} projetado --- exatamente a função do diodo de roda
livre, aqui desempenhada por um resistor.}
\end{questao}

\blocodesafio

\begin{questao}[4]{}
\textbf{(Dedução --- resposta do circuito RL.)} Um circuito RL série com fonte
$V$ tem a chave fechada em $t=0$, com $i(0)=0$.
\begin{itensq}
\item Escreva a equação diferencial e resolva-a por separação de variáveis.
\item Identifique $\tau$ e o valor final.
\item Obtenha $v_L(t)$ e verifique a LKT em $t=0^{+}$ e $t\to\infty$.
\end{itensq}
\resposta{$i(t)=\dfrac{V}{R}\left(1-e^{-Rt/L}\right)$;
$\tau=L/R$; $v_L=Ve^{-Rt/L}$}
\resolucao{(a) Pela LKT,
\[
V=Ri+L\frac{\mathrm{d}i}{\mathrm{d}t}
\;\Longrightarrow\;
\frac{\mathrm{d}i}{V/R-i}=\frac{R}{L}\,\mathrm{d}t.
\]
Integrando de $0$ a $t$ (com $i(0)=0$):
\[
-\ln\!\left(\frac{V/R-i}{V/R}\right)=\frac{R}{L}t
\;\Longrightarrow\;
i(t)=\frac{V}{R}\left(1-e^{-Rt/L}\right).
\]
(b) O expoente deve ser adimensional, o que identifica
$\tau=L/R$; e $i(\infty)=V/R$, coerente com o indutor tornando-se um curto.\\
(c)
\[
v_L=L\frac{\mathrm{d}i}{\mathrm{d}t}
=L\cdot\frac{V}{R}\cdot\frac{R}{L}e^{-Rt/L}=V\,e^{-Rt/L}.
\]
\emph{Em $t=0^{+}$:} $v_L=V$ e $Ri=0$; soma $=V$ \checkmark.\\
\emph{Em $t\to\infty$:} $v_L=0$ e $Ri=V$; soma $=V$ \checkmark.\\
\textbf{Estrutura geral:} a solução é (regime permanente) $+$ (transitório
exponencial). Essa decomposição vale para \emph{qualquer} circuito de primeira
ordem e dispensa refazer a integração a cada problema.}
\end{questao}

\begin{questao}[4]{}
\textbf{(Integração --- energia na descarga.)} Um indutor com corrente inicial
$I_0$ é descarregado através de $R$.
\begin{itensq}
\item Escreva $i(t)$ e a potência instantânea no resistor.
\item Integre para obter a energia total dissipada.
\item Mostre que o resultado independe de $R$ e interprete.
\end{itensq}
\resposta{$W=\frac12LI_0^{2}$, independente de $R$}
\resolucao{(a) Sem fonte, $Ri+L\,\mathrm{d}i/\mathrm{d}t=0$, cuja solução é
$i(t)=I_0e^{-t/\tau}$ com $\tau=L/R$. Então
\[
p_R(t)=Ri^{2}=RI_0^{2}e^{-2t/\tau}.
\]
(b)
\[
W=\int_0^{\infty}RI_0^{2}e^{-2t/\tau}\,\mathrm{d}t
=RI_0^{2}\left[-\frac{\tau}{2}e^{-2t/\tau}\right]_0^{\infty}
=\frac{RI_0^{2}\tau}{2}.
\]
Substituindo $\tau=L/R$:
\[
W=\frac{RI_0^{2}}{2}\cdot\frac{L}{R}=\frac12LI_0^{2}.
\]
(c) O $R$ \textbf{cancela} exatamente: ele aparece na potência
($\propto R$) e desaparece na duração ($\tau\propto1/R$). Aumentar $R$ dez
vezes decuplica a potência inicial e divide o tempo por dez --- a área sob a
curva não muda.\\
\textbf{Significado:} a energia dissipada é determinada pelo \emph{estado
inicial} do indutor, não pelo caminho de descarga. É a mesma lógica que faz um
capacitor descarregado entregar sempre $\frac12CV_0^{2}$, e confirma por
integração o argumento de conservação usado no bloco anterior.}
\end{questao}

\begin{questao}[4]{}
\textbf{(Projeto de proteção.)} Uma bobina de relé com $L=\SI{0.2}{\henry}$ e
$r=\SI{100}{\ohm}$ opera em $\SI{24}{\volt}$. O transistor de comando suporta
no máximo $\SI{60}{\volt}$.
\begin{itensq}
\item Determine a corrente de operação e a energia armazenada.
\item Dimensione um grampeador resistivo (diodo em série com $R_g$) que
      respeite o limite do transistor.
\item Compare o tempo de extinção com o do diodo simples.
\end{itensq}
\resposta{(a) $\SI{0.24}{\ampere}$, $\SI{5.76}{\milli\joule}$;
(b) $R_g\le\SI{150}{\ohm}$; (c) $\SI{0.8}{\milli\second}$ contra
$\SI{2}{\milli\second}$}
\resolucao{(a) $I=24/100=\SI{0.24}{\ampere}$;
$W=\frac12(0{,}2)(0{,}24)^{2}=\SI{5.76}{\milli\joule}$.\\
(b) Na abertura, a corrente $\SI{0.24}{\ampere}$ é forçada pelo diodo e por
$R_g$. A tensão no coletor sobe para $V+R_gI$ (desprezando o diodo). Exigindo
$\le\SI{60}{\volt}$:
\[
24+R_g(0{,}24)\le60
\;\Longrightarrow\;
R_g\le\frac{36}{0{,}24}=\SI{150}{\ohm}.
\]
Adote $R_g=\SI{150}{\ohm}$ e verifique a dissipação: o pulso entrega
$\SI{5.76}{\milli\joule}$ e, a $\SI{10}{\hertz}$ de comutação, isso são
$\SI{57.6}{\milli\watt}$ médios --- um resistor de
$\frac14\,\si{\watt}$ basta, desde que suporte o pico de
$0{,}24^{2}(150)=\SI{8.6}{\watt}$ instantâneos.\\
(c) \emph{Com grampeador:} $\tau=L/(r+R_g)=0{,}2/250=\SI{0.8}{\milli\second}$.\\
\emph{Com diodo simples:} $\tau=L/r=0{,}2/100=\SI{2}{\milli\second}$.\\
\textbf{Compromisso quantificado:} o grampeador reduz o tempo de desarme em
$60\%$, ao custo de submeter o transistor a $\SI{60}{\volt}$ em vez de
$\SI{24.7}{\volt}$. Quanto maior $R_g$, mais rápido o relé desarma e maior a
tensão no interruptor --- o projeto consiste em usar toda a margem do
transistor, e nada além dela.}
\end{questao}

\begin{questao}[4]{}
\textbf{(Projeto de armazenamento.)} Deseja-se um indutor que armazene
$\SI{5}{\joule}$ conduzindo $\SI{10}{\ampere}$.
\begin{itensq}
\item Determine $L$.
\item Para um núcleo toroidal com $A=\SI{10}{\centi\meter\squared}$,
      $\ell=\SI{0.25}{\meter}$ e $\mu_r=2000$, determine o número de espiras.
\item Verifique a densidade de fluxo e avalie a viabilidade.
\end{itensq}
\resposta{(a) $L=\SI{0.1}{\henry}$; (b) $N\approx35$ espiras;
(c) $B\approx\SI{3.5}{\tesla}$ --- \textbf{inviável}, o núcleo satura}
\resolucao{(a) $W=\frac12LI^{2}\Rightarrow
L=\dfrac{2W}{I^{2}}=\dfrac{10}{100}=\SI{0.1}{\henry}$.\\
(b) De $L=\dfrac{\mu_r\mu_0N^{2}A}{\ell}$:
\[
N=\sqrt{\frac{L\ell}{\mu_r\mu_0A}}
=\sqrt{\frac{0{,}1\cdot0{,}25}{2000(4\pi\times10^{-7})(10^{-3})}}
=\sqrt{9{,}95}\approx35\ \text{espiras}.
\]
(c) O fluxo é
\[
B=\frac{\mu_r\mu_0NI}{\ell}
=\frac{2000(4\pi\times10^{-7})(35)(10)}{0{,}25}\approx\SI{3.5}{\tesla}.
\]
Materiais ferromagnéticos comuns saturam entre $1{,}2$ e
$\SI{1.8}{\tesla}$. O projeto é \textbf{inviável}: muito antes de atingir
$\SI{10}{\ampere}$, o núcleo satura, $\mu_r$ despenca e $L$ some.\\
\textbf{Diagnóstico e correção.} A energia armazenada em material
ferromagnético é limitada justamente porque $\mu_r$ alto significa $B$ alto para
pouca corrente. A solução clássica é introduzir um \textbf{entreferro}: ele
reduz a permeabilidade efetiva (exigindo mais espiras), mas permite corrente
muito maior antes da saturação --- e a maior parte da energia passa a residir
no entreferro, onde $u=B^{2}/2\mu_0$ é enorme.\\
\textbf{Lição:} calcular $L$ e $N$ não basta. Sem verificar $B$, obtém-se um
projeto aritmeticamente correto e fisicamente impossível.}
\end{questao}

\begin{questao}[4]{}
\textbf{(Projeto de constante de tempo.)} Deseja-se $\tau=\SI{20}{\milli\second}$
com $R=\SI{50}{\ohm}$.
\begin{itensq}
\item Determine $L$.
\item Avalie a viabilidade física dessa indutância.
\item Proponha duas alternativas para obter o mesmo $\tau$ de forma prática.
\end{itensq}
\resposta{(a) $L=\SI{1}{\henry}$; (c) aumentar $R$, ou usar RC equivalente}
\resolucao{(a) $L=\tau R=0{,}020(50)=\SI{1}{\henry}$.\\
(b) \textbf{Um henry é um componente grande.} Exige núcleo ferromagnético e
centenas de espiras; a resistência do próprio enrolamento facilmente atinge
dezenas de ohms --- comparável aos $\SI{50}{\ohm}$ do projeto, o que
\emph{altera} o $\tau$ pretendido. Some-se a isso massa, custo, sensibilidade a
campos externos e saturação.\\
(c) \textbf{Alternativa 1 --- elevar $R$.} Com $R=\SI{500}{\ohm}$,
$L=\SI{10}{\henry}$... pior. O caminho correto é o inverso: fixar um $L$
razoável e ajustar $R$. Com $L=\SI{10}{\milli\henry}$, obter
$\tau=\SI{20}{\milli\second}$ exigiria $R=\SI{0.5}{\ohm}$ --- baixo demais para
ser prático em circuitos de sinal.\\
\textbf{Alternativa 2 --- usar RC.} Um capacitor de
$\SI{20}{\micro\farad}$ com $\SI{1}{\kilo\ohm}$ dá $\tau=\SI{20}{\milli\second}$
em um componente barato, pequeno e estável.\\
\textbf{Conclusão de engenharia:} para constantes de tempo da ordem de
milissegundos ou mais, \textbf{RC vence RL} quase sempre. Indutores são
preferíveis quando se precisa de armazenamento de energia com corrente elevada
ou de comportamento em frequência --- não para gerar atrasos.}
\end{questao}

\begin{questao}[4]{}
\textbf{(Demonstração --- acoplamento em série.)} Duas bobinas acopladas em
série conduzem a mesma corrente $i$.
\begin{itensq}
\item Deduza $L_{eq}=L_1+L_2\pm2M$ a partir das tensões induzidas.
\item Demonstre que $M\le\sqrt{L_1L_2}$, usando o fato de a energia ser não
      negativa.
\item Explique o significado da convenção dos pontos.
\end{itensq}
\resposta{(b) da condição $W\ge0$ para toda corrente segue
$M^{2}\le L_1L_2$}
\resolucao{(a) Cada bobina sofre autoindução \emph{e} indução mútua:
\[
v_1=L_1\frac{\mathrm{d}i}{\mathrm{d}t}\pm M\frac{\mathrm{d}i}{\mathrm{d}t},
\qquad
v_2=L_2\frac{\mathrm{d}i}{\mathrm{d}t}\pm M\frac{\mathrm{d}i}{\mathrm{d}t}.
\]
Somando e comparando com
$v=L_{eq}\,\mathrm{d}i/\mathrm{d}t$:
\[
L_{eq}=L_1+L_2\pm2M.
\]
(b) Com correntes independentes $i_1$ e $i_2$, a energia do par é
\[
W=\frac12L_1i_1^{2}+\frac12L_2i_2^{2}+Mi_1i_2.
\]
Como um sistema passivo não pode armazenar energia negativa, $W\ge0$ para
\emph{quaisquer} $i_1,i_2$. Tomando $i_1=x$ e $i_2=-1$:
\[
\frac12L_1x^{2}-Mx+\frac12L_2\ge0\quad\forall x,
\]
o que exige discriminante não positivo:
\[
M^{2}-L_1L_2\le0
\;\Longrightarrow\;
M\le\sqrt{L_1L_2}
\;\Longrightarrow\;
k=\frac{M}{\sqrt{L_1L_2}}\le1.
\]
(c) Os pontos marcam terminais \textbf{magneticamente correspondentes}:
correntes que entram pelos terminais pontuados produzem fluxos que se
\emph{reforçam}. Se ambas as correntes entram (ou ambas saem) pelos pontos, o
termo mútuo é positivo; caso contrário, negativo.\\
\textbf{Por que a convenção é necessária:} o esquema elétrico não informa o
sentido de enrolamento das bobinas, que é uma propriedade \emph{física} da
construção. Sem os pontos, $L_{eq}$ seria ambíguo entre
$\SI{19}{\milli\henry}$ e $\SI{7}{\milli\henry}$ --- e, em um transformador, a
polaridade da saída ficaria indeterminada.}
\end{questao}

\begin{questao}[4]{}
\textbf{(Indutor real em regime.)} Um indutor com $L=\SI{0.5}{\henry}$ e
$r=\SI{2}{\ohm}$ é alimentado por $\SI{10}{\volt}$ CC.
\begin{itensq}
\item Determine a corrente, a energia armazenada e a potência dissipada em
      regime.
\item Determine o tempo necessário para atingir $99\%$ da corrente final.
\item Discuta o "custo de manutenção" da energia armazenada.
\end{itensq}
\resposta{(a) $\SI{5}{\ampere}$, $\SI{6.25}{\joule}$, $\SI{50}{\watt}$;
(b) $\approx\SI{1.15}{\second}$}
\resolucao{(a) $I=10/2=\SI{5}{\ampere}$;
$W=\frac12(0{,}5)(25)=\SI{6.25}{\joule}$;
$P=rI^{2}=2(25)=\SI{50}{\watt}$.\\
(b) $\tau=L/r=0{,}5/2=\SI{0.25}{\second}$, e $99\%$ correspondem a
$t=\tau\ln100=0{,}25(4{,}605)=\SI{1.15}{\second}$.\\
(c) \textbf{O contraste é gritante.} Carregar o campo custou
$\SI{6.25}{\joule}$; mantê-lo custa $\SI{50}{\watt}$ \emph{continuamente}.
Em apenas $\SI{0.125}{\second}$ de operação, a dissipação já iguala toda a
energia armazenada; em um minuto, dissipam-se $\SI{3000}{\joule}$ --- $480$
vezes o conteúdo do campo.\\
\textbf{Figura de mérito:} a razão $W/P=L/(2r)=\tau/2$ tem dimensão de tempo e
mede a "qualidade" do indutor em CC. Aqui vale $\SI{0.125}{\second}$.\\
\textbf{Consequência:} indutores são péssimos \emph{reservatórios} de energia
para armazenamento prolongado --- ao contrário de baterias ou capacitores, que
não exigem corrente para manter o estado. Eles são excelentes para
\emph{transferir} energia em ciclos curtos, que é exatamente o que fazem em
conversores chaveados, onde o ciclo dura microssegundos.}
\end{questao}

\begin{questao}[4]{}
\textbf{(Conversor chaveado.)} Um conversor opera a $\SI{50}{\kilo\hertz}$ com
razão cíclica $D=0{,}5$, aplicando $\SI{12}{\volt}$ sobre um indutor de
$\SI{100}{\micro\henry}$ durante a condução.
\begin{itensq}
\item Determine a ondulação de corrente pico a pico.
\item Determine a corrente de pico se a componente contínua for
      $\SI{3}{\ampere}$.
\item O indutor satura em $\SI{3.5}{\ampere}$. Proponha e quantifique duas
      correções.
\end{itensq}
\resposta{(a) $\Delta i=\SI{1.2}{\ampere}$; (b) $I_{pico}=\SI{3.6}{\ampere}$;
(c) satura --- elevar $L$ ou a frequência}
\resolucao{(a) Durante o tempo de condução
$t_{on}=D/f=0{,}5/\pot{5}{4}=\SI{10}{\micro\second}$:
\[
\Delta i=\frac{V\,t_{on}}{L}=\frac{12(10^{-5})}{10^{-4}}=\SI{1.2}{\ampere}.
\]
(b) A ondulação se distribui simetricamente em torno da média:
\[
I_{pico}=3+\frac{1{,}2}{2}=\SI{3.6}{\ampere}\ >\ \SI{3.5}{\ampere}.
\]
\textbf{O indutor satura} --- e note que a corrente \emph{média}
($\SI{3}{\ampere}$) estava confortavelmente abaixo do limite. Dimensionar pela
média é o erro clássico.\\
(c) Como $\Delta i=\dfrac{VD}{fL}$, há dois caminhos:\\
\textbf{(i) Aumentar $L$.} Para $I_{pico}\le3{,}5$ é preciso
$\Delta i\le\SI{1}{\ampere}$, ou seja
\[
L\ge\frac{12(10^{-5})}{1}=\SI{120}{\micro\henry}.
\]
Custo: indutor maior e mais caro.\\
\textbf{(ii) Aumentar a frequência.} Mantendo
$L=\SI{100}{\micro\henry}$, exige-se
\[
f\ge\frac{12(0{,}5)}{10^{-4}(1)}=\SI{60}{\kilo\hertz}.
\]
Custo: mais perdas de comutação no transistor e maior interferência
eletromagnética.\\
\textbf{Compromisso:} $\Delta i\propto1/(fL)$ --- o produto $fL$ é o que
importa. Elevar a frequência permite indutores menores, e é exatamente essa a
tendência que tornou as fontes chaveadas modernas tão compactas.}
\end{questao}

\begin{questao}[4]{}
\textbf{(Demonstração --- continuidade da corrente.)} Prove que a corrente em
um indutor ideal não pode sofrer descontinuidade, usando um argumento
energético.
\begin{itensq}
\item Suponha um salto de $I_1$ para $I_2$ em intervalo $\Delta t\to0$ e
      analise a tensão.
\item Analise a potência e a energia envolvidas.
\item Explique por que capacitores em paralelo apresentam o problema dual.
\end{itensq}
\resposta{Um salto exigiria tensão e potência infinitas; a energia é finita,
mas a taxa de transferência não pode ser}
\resolucao{(a) Um salto implica
\[
v=L\lim_{\Delta t\to0}\frac{I_2-I_1}{\Delta t}\to\infty.
\]
Nenhuma fonte real fornece tensão infinita.\\
(b) A potência $p=vi$ também diverge. A \emph{energia} necessária é finita
---
\[
\Delta W=\frac12L\left(I_2^{2}-I_1^{2}\right),
\]
--- mas ela teria de ser transferida em tempo nulo, exigindo potência
infinita. \textbf{É a taxa, não a quantidade, que torna o salto impossível.}\\
(c) \textbf{Problema dual:} dois capacitores com tensões diferentes ligados em
paralelo exigiriam salto de \emph{tensão}, o que demandaria corrente infinita
($i=C\,\mathrm{d}v/\mathrm{d}t$). Na prática, a resistência inevitável dos
condutores limita a corrente e dissipa a diferença de energia --- e demonstra-se
que a perda independe do valor dessa resistência, resultado análogo ao da
descarga do indutor.\\
\textbf{Regra prática:} nunca ligue em série indutores com correntes distintas,
nem em paralelo capacitores com tensões distintas. Em ambos os casos, o modelo
ideal falha e o circuito real encontra uma saída --- arco elétrico no primeiro
caso, pico de corrente no segundo.}
\end{questao}

\begin{questao}[4]{}
\textbf{(Modelagem --- sensor indutivo.)} Um sensor de proximidade usa uma
bobina cuja indutância varia com a distância $x$ a um alvo ferromagnético,
segundo $L(x)=L_0\left(1+\dfrac{a}{x}\right)$, com
$L_0=\SI{10}{\milli\henry}$ e $a=\SI{1}{\milli\meter}$.
\begin{itensq}
\item Determine $L$ para $x=\SI{1}{\milli\meter}$ e $x=\SI{5}{\milli\meter}$.
\item Obtenha a sensibilidade $\mathrm{d}L/\mathrm{d}x$ e avalie nos dois
      pontos.
\item Discuta a consequência para o projeto do sensor.
\end{itensq}
\resposta{(a) $\SI{20}{\milli\henry}$ e $\SI{12}{\milli\henry}$;
(b) $-\SI{10}{\milli\henry\per\milli\meter}$ e
$-\SI{0.4}{\milli\henry\per\milli\meter}$}
\resolucao{(a)
\[
L(1)=10(1+1)=\SI{20}{\milli\henry};
\qquad
L(5)=10(1+0{,}2)=\SI{12}{\milli\henry}.
\]
(b)
\[
\frac{\mathrm{d}L}{\mathrm{d}x}=-\frac{L_0a}{x^{2}}.
\]
Em $x=\SI{1}{\milli\meter}$:
$-10(1)/1=-\SI{10}{\milli\henry\per\milli\meter}$.\\
Em $x=\SI{5}{\milli\meter}$:
$-10(1)/25=-\SI{0.4}{\milli\henry\per\milli\meter}$.\\
A sensibilidade cai com $1/x^{2}$: é \textbf{25 vezes menor} a
$\SI{5}{\milli\meter}$.\\
(c) \textbf{Três consequências.} \emph{(i)} O sensor é fortemente
\textbf{não linear} --- a leitura bruta de $L$ não é proporcional a $x$, e a
conversão exige linearização ou tabela. \emph{(ii)} A \textbf{faixa útil} é
curta: além de poucos milímetros, a variação se confunde com a deriva térmica
da própria bobina. \emph{(iii)} A \textbf{resolução} depende da posição: perto
do alvo, décimos de micrômetro são detectáveis; longe, nem décimos de
milímetro.\\
\textbf{Solução usual:} operar sempre na região próxima e usar o sensor como
detector de \emph{limiar} (presença/ausência), em vez de medidor de distância
--- que é exatamente como os sensores indutivos industriais são
especificados.}
\end{questao}
